\subsection{Code Generation}
\noindent \textbf{Language Models for Code Generation.} 
Starting with Codex~\citep{chen2021evaluating}, there are over a dozen code \llms{}. These include CodeT5~\citep{wang2021codet5, wang2023codet5+}, CodeGen~\citep{nijkamp2022codegen}, SantaCoder~\citep{allal2023santacoder}, StarCoder~\citep{li2023starcoder}, AlphaCode~\citep{li2022competition}, InCoder~\citep{fried2022incoder}, and CodeGeeX~\citep{zheng2023codegeex}. 
As of May 2024, \llamabase{} and \deepseek{}~\citep{bi2024deepseek}, \starcoder{}~\cite{starcoder2,li2023starcoder} and \cllama{}~\citep{roziere2023code} are the most popular open models. 
Many downstream models resulted from fine-tuning them on synthetically generated data, such as \wizardlong{} \citep{luo2023wizardcoder}, \magicoders{} \citep{wei2023magicoder}, and \phind{}.

\vspace{1pt}
\noindent \textbf{Code Generation Benchmarks.} 
Many benchmarks have been proposed to compare and evaluate these models. 
These primarily focus on natural language to Python code generation: 
\humaneval{} \citep{chen2021evaluating}, 
\humanevalplus{} \citep{liu2023your}, 
\apps{} \citep{hendrycks2021measuring}, 
\contests{} \citep{li2022competition}, 
\mbpp{} \citep{austin2021program}, 
L2CEval \citep{ni2023l2ceval}. 
Their variants have been proposed to cover more languages, \citep{wang2022recode, zheng2023codegeex, cassano2022multipl, athiwaratkun2022multi}.
Many benchmarks have focused on code generation in APIs. Benchmarks like DS-1000 \citep{lai2023ds}, ARCADE \citep{yin2022natural}, NumpyEval \citep{zhang2023toolcoder}, and PandasEval \citep{jain2022jigsaw} focus on data science APIs. 
Other benchmarks measure using broader APIs or 
general software engineering tasks, such as JuICe~\citep{agashe2019juice}, APIBench~\citep{patil2023gorilla}, RepoBench~\citep{liu2023repobench}, ODEX~\citep{wang2022execution}, SWE-Bench~\citep{jimenez2023swe}, GoogleCodeRepo~\citep{shrivastava2023repository}, RepoEval~\citep{zhang2023repocoder}, and Cocomic-Data~\citep{ding2022cocomic}.~

A few benchmarks specifically measure competitive programming, such as 
APPS \citep{hendrycks2021measuring}, 
CodeContests \citep{li2022competition}, 
CodeScope \citep{yan2023codescope},
xCodeEval \citep{khan2023xcodeeval},
and LeetCode-Hard \citep{shinn2023reflexion}, and TACO \citep{li2023taco}. 
Methods such as AlphaCode \citep{li2022competition}, AlphaCode 2\citep{team2023gemini}, ALGO \citep{zhang2023algo}, Parsel \citep{zelikman2022parsel}, code cleaning \citep{jain2023llm}, code explainations \citep{li2023explaining}, analogical reasoning \citep{yasunaga2023large}, and AlphaCodium \citep{ridnik2024code} have been pushing the boundaries of what is possible with \llms{} in this domain. The biggest differentiating factor between \livecodebench{} and these benchmarks is that our benchmark is \textbf{continuously updated}, \textbf{problem curation with balanced difficulty}, \textbf{higher tests and problem quality}, and contains \textbf{more scenarios} such as code repair, code execution, and test output prediction capturing more facets for building agentic coding systems.

%\redo{TODO are these two citations still needed?}
%\cite{li2023explaining}
%\cite{}
\subsection{Holistic Tasks}
\livecodebench{} considers self-repair, test output prediction, and code execution as additional scenarios. Below we note pertinent related work for these domains. 

%Finally, there are a variety of benchmarks for other tasks, such as code translation \citep{roziere2020unsupervised, zhu2022xlcost, ahmad-etal-2021-avatar}, test case generation \citep{tufano2022methods2test, watson2020learning}, code search \citep{husain2019codesearchnet}, type prediction \citep{mir2022type4py, wei2023typet5, malik2019nl2type}, commit message generation \citep{liu2020atom}, code summarization \citep{leclair2019neural, iyer2016summarizing, barone2017parallel, hasan2021codesc, alon2018code2seq}, code security \citep{liguori2022can, pearce2022asleep, tony2023llmseceval}, program repair \citep{jiang2023impact, xia2022practical, tufano2019empirical, haque2022fixeval, jin2023inferfix, gupta2017deepfix, berabi2021tfix}, performance optimization \citep{garg2022deepperf, madaan2023learning}, and so on.

\vspace{1pt}
\noindent \textbf{Code Repair.} 
\citep{chen2023teaching, olausson2023demystifying, madaan2023self, peng2023check, zhang2023self} have investigated self-repair for existing code \llm{} benchmarks. Particularly, these methods use error feedback for models to improve inspiring our code repair scenario.

\vspace{1pt}
\noindent \textbf{Code Execution.} 
Code execution was first studied in \citet{austin2021program, nye2021show}
\livecodebench{}'s execution scenario is particularly inspired by CRUXEval \citep{gu2024cruxeval}, a recent benchmark measuring the reasoning and execution abilities of code \llms{}. 
We differ from CRUXEval in that our benchmark is live, and our functions are more complex and human-produced (unlike Code Llama generations in CRUXEval).% and are usually more complex.
%There are a few differences between our code execution scenario and CRUXEval, allowing our execution task to provide a complementary perspective. First, our benchmark is live, which will defend against contamination.
%The functions in our benchmark are written by human programmers and are natural, while functions in CRUXEval are distilled by Code Llama. 
%Third, functions in our benchmark have a natural language intent, while functions in CRUXEval are more arbitrary. 
%Third, our functions tend to be more complex than those in CRUXEval due to the nature of their source.

\vspace{1pt}
\noindent \textbf{Test Generation.} 
Test generation using \llms{} has been explored in \citep{yuan2023no, 10329992, tufano2022methods2test, watson2020learning}. 
Furthermore, \citet{chen2022codet} demonstrated that \llms{} can assist in generating test case inputs/outputs for competitive programming problems, thereby improving the accuracy of the generated code, thus inspiring our test generation scenario.
However, \livecodebench{}'s test generation scenario is unique in that it decouples the test inputs and outputs allowing more proper evaluations.

Finally, some works have additionally studied other tasks and scenarios like type prediction \citep{mir2022type4py, wei2023typet5, malik2019nl2type}, code summarization \citep{leclair2019neural, iyer2016summarizing, barone2017parallel, hasan2021codesc, alon2018code2seq}, code security \citep{liguori2022can, pearce2022asleep, tony2023llmseceval}, etc. 

%\vspace{1pt}
%\noindent \textbf{Contamination}:
\subsection{Contamination}
Data contamination and test-case leakage have received considerable attention ~\cite{yonatan2024proving,golchin2023time,weller2023according,roberts2024to} as \llms{} might be getting trained on benchmarks. 
~\citet{sainz2023chatgpt} demonstrated contamination by simply prompting the model to highlight its contamination. 
Some detection methods have also been built to avoid these cases~\citep{shi2023detecting,zhou2023don}. For code, \citet{riddell2024quantifying} use edit distance and AST-based semantic-similarity to detect contamination. 
%Our work addresses it by live updates.% such that people can build new tests, models, metrics, and prompting strategies. 